\section{问题三模型建立和求解}

\subsection{三弹联合时序模型}

设 FY1 的共享航向角和速度为 $\theta$ 与 $v$，第 $i$ 枚烟幕弹的投放时刻、引信延时和起爆时刻分别为 $t_{r,i}$、$\tau_i$ 与 $t_{e,i}=t_{r,i}+\tau_i$。用非负松弛量 $s_{12},s_{23}$ 编码投放间隔：
\begin{equation}
 t_{r,2}=t_{r,1}+1+s_{12},\qquad
 t_{r,3}=t_{r,2}+1+s_{23},\qquad s_{12},s_{23}\ge0.
\end{equation}
第 $i$ 枚烟幕弹的投放点和起爆点由同一条航迹确定：
\begin{align}
 \bm r_{r,i}&=\bm u_0+v\bm e(\theta)t_{r,i},\\
 \bm r_{e,i}&=\bm u_0+v\bm e(\theta)t_{e,i}
 +(0,0,-\tfrac12g\tau_i^2)^{\mathsf T}.
\end{align}
分别记三枚烟幕弹的有效集合为 $\Omega_i$。总时长必须取并集 $\Omega=\bigcup_i\Omega_i$ 的测度，不能直接累加单弹时长，否则重叠部分会被重复计算。

\subsection{优化模型汇总}

问题三的八维优化模型为
\begin{equation}
\begin{aligned}
&\max_{\bm y}\quad T_{\mathrm{eff}}^{(3)}(\bm y)
=\mu\!\left(\bigcup_{i=1}^{3}\Omega_i(\bm y)\right),\\[-1mm]
&\text{\rm s.t.}\quad
\left\{\begin{aligned}
&\bm y=(\theta,v,t_{r,1},s_{12},s_{23},\tau_1,\tau_2,\tau_3)^{\mathsf T},\\
&0\le\theta<2\pi,\qquad 70\le v\le140,\\
&t_{r,1}\ge0,\quad s_{12},s_{23}\ge0,\\
&t_{r,2}=t_{r,1}+1+s_{12},\quad t_{r,3}=t_{r,2}+1+s_{23},\\
&\tau_i\ge0,\qquad t_{e,i}=t_{r,i}+\tau_i\le t_M,\\
&\bm r_{r,i}=\bm u_0+v\bm e(\theta)t_{r,i},\\
&\bm r_{e,i}=\bm u_0+v\bm e(\theta)t_{e,i}+(0,0,-\tfrac12g\tau_i^2)^{\mathsf T},\\
&z_{e,i}=1800-\tfrac12g\tau_i^2\ge0,\\
&\bm c_i(t)=\bm r_{e,i}-(0,0,3(t-t_{e,i}))^{\mathsf T},\\
&t_{e,i}\le t\le\min\{t_{e,i}+20,t_M\},\\
&D_i(t)=\displaystyle\max_{\bm p\in\mathcal T}d_i(\bm p,t),\\
&\Omega_i=\{t:D_i(t)\le10\},\qquad i=1,2,3.
\end{aligned}\right.
\end{aligned}
\label{eq:q3-complete-model}
\end{equation}
该模型同时约束共享航迹、投放间隔、自由落体高度损失、烟幕寿命和完整圆柱遮蔽几何。由松弛量编码产生的候选解天然满足最小投放间隔，可减少惩罚函数造成的无效搜索。

\subsection{联合优化与精确复核}

真实目标由多个阈值事件和区间并运算构成，具有大面积零值区、多个航向盆地和窄时间窗。搜索阶段先用平滑遮蔽概率近似引导差分进化，再在全航向、外向和内向扇区执行多初值搜索；同时加入首尾时间窗拼接构造，避免种群遗漏靠近区间衔接边界的高质量策略。差分进化适合处理不可导、多峰连续优化问题\cite{storn1997de}。

候选排序使用真实几何判据：计算每枚烟幕球到完整圆柱各条视线段的最大距离，以 Brent 法求解 $D_i(t)=10$ 的进出事件\cite{brent1973algorithms}，再合并全部有效区间。最终以每个圆周 $5760$ 个角向点和 $0.005\,\mathrm{s}$ 的扫描步长复核。流程见\cref{fig:q3-flow}。

\begin{figure}[!htbp]
 \centering
 \includegraphics[width=0.72\textwidth]{fig_flow_q3.pdf}
 \caption{问题三三弹联合优化与区间并集复核流程}
 \label{fig:q3-flow}
\end{figure}

\subsection{投放策略与遮蔽结果}

最优候选取航向角 $179.647479^{\circ}$、速度 $139.998280\,\mathrm{m/s}$，对应速度向量为
\begin{equation}
 \bm v_U=(-139.995630,0.861355,0)^{\mathsf T}\,\mathrm{m/s}.
\end{equation}
三次投放分别发生在 $0.003$、$3.702$ 和 $5.569\,\mathrm{s}$，相邻间隔为 $3.699$ 和 $1.867\,\mathrm{s}$。具体策略见\cref{tab:q3-strategy}。

\begin{table}[!htbp]
 \centering
 \caption{问题三 FY1 三枚烟幕弹投放策略}
 \label{tab:q3-strategy}
 \resizebox{\textwidth}{!}{
 \begin{tabular}{cccccc}
  \toprule[1.5pt]
  编号 & 投放时刻/s & 投放点/m & 起爆时刻/s & 起爆点/m & 单弹时长/s \\
  \midrule[1pt]
  1 & $0.003$ & $(17799.580,0.003,1800.000)$ & $3.614$ & $(17294.056,3.113,1736.107)$ & $3.834877$ \\
  2 & $3.702$ & $(17281.736,3.189,1800.000)$ & $9.039$ & $(16534.579,7.786,1660.431)$ & $2.591030$ \\
  3 & $5.569$ & $(17020.364,4.797,1800.000)$ & $11.610$ & $(16174.651,10.000,1621.181)$ & $1.204687$ \\
  \bottomrule[1.5pt]
 \end{tabular}}
\end{table}

三枚弹的有效区间分别为
\begin{equation}
 [5.205021,9.039899],\quad[9.039000,11.630030],\quad[11.610000,12.814687].
\end{equation}
相邻区间在端点附近交叠，因而拼接成
\begin{equation}
 \Omega=[5.205021,12.814687],\qquad
 T_{\mathrm{eff}}^{(3)}=7.609665451\,\mathrm{s}\approx\boxed{7.610\,\mathrm{s}}.
\end{equation}
与早投边界候选相比，新方案通过延迟起爆把三个云团依次布置在 M1 视线前方，第三枚弹获得正边际收益并延长尾部窗口。该方案同步写入 \texttt{result1.xlsx}，并作为问题五中 FY1 对 M1 的子策略。

\begin{figure}[!htbp]
 \centering
 \includegraphics[width=0.96\textwidth]{q3_optimization_diagnostics.pdf}
 \caption{问题三多盆地搜索与三弹最坏视线距离}
 \label{fig:q3-diagnostics}
\end{figure}

\subsection{收敛性与稳健性检验}

角向采样从 $1440$ 点加密到 $5760$ 点时，联合时长变化小于 $10^{-4}\,\mathrm{s}$。三个起爆高度分别为 $1736.107$、$1660.431$ 和 $1621.181\,\mathrm{m}$，均远高于地面；全部有效区间都处于相应烟幕起爆后的 $20\,\mathrm{s}$ 内，并早于 M1 到达假目标。

最优区间的衔接裕量较小，说明时序误差比位置小扰动更敏感。本文通过事件求根而非固定网格计数确定端点，避免把端点接触误判为断裂。不同随机种子可得到参数略有差异但时长相近的候选，因此本文将其表述为经多盆地搜索和高密度复核的最优候选，不声称解析意义上的全局最优。
